% ====================================================================
% EXEMPLO-FORMULARIO.TEX — COMO CRIAR CAMPOS, CAMPO LONGO, CAIXAS E TABELA
% ====================================================================
% Este arquivo é só um EXEMPLO de uso, não é chamado pelo main.tex.
% Para usar algo assim de verdade no corpo do documento, dê
% % ====================================================================
% EXEMPLO-FORMULARIO.TEX — COMO CRIAR CAMPOS, CAMPO LONGO, CAIXAS E TABELA
% ====================================================================
% Este arquivo é só um EXEMPLO de uso, não é chamado pelo main.tex.
% Para usar algo assim de verdade no corpo do documento, dê
% % ====================================================================
% EXEMPLO-FORMULARIO.TEX — COMO CRIAR CAMPOS, CAMPO LONGO, CAIXAS E TABELA
% ====================================================================
% Este arquivo é só um EXEMPLO de uso, não é chamado pelo main.tex.
% Para usar algo assim de verdade no corpo do documento, dê
% % ====================================================================
% EXEMPLO-FORMULARIO.TEX — COMO CRIAR CAMPOS, CAMPO LONGO, CAIXAS E TABELA
% ====================================================================
% Este arquivo é só um EXEMPLO de uso, não é chamado pelo main.tex.
% Para usar algo assim de verdade no corpo do documento, dê
% \input{secoes/exemplo-formulario.tex} dentro de main.tex (onde hoje
% está \input{secoes/teste.tex}).
%
% As macros usadas abaixo (\campotexto, \campolongo, \caixa,
% \campotabela, \checktabela) são definidas em
% estrut/campos-formulario.tex — não precisa mexer lá, só usá-las.
%
% REGRA OBRIGATÓRIA: todo campo só funciona dentro do ambiente Form,
% como abaixo. NOMES (1º argumento de cada macro) DEVEM SER ÚNICOS
% no documento inteiro.
% ====================================================================

\begin{Form}

\section{Exemplo de Formulário}

% --------------------------------------------------------------
% 1. Campo de texto simples (uma linha)
% --------------------------------------------------------------
Nome do responsável: \campotexto{exemplo_nome}{6cm}

\vspace{0.5cm}

% --------------------------------------------------------------
% 2. Campo de texto longo (múltiplas linhas)
% --------------------------------------------------------------
Observações gerais:

\campolongo{exemplo_observacoes}{3cm}

\vspace{0.5cm}

% --------------------------------------------------------------
% 3. Caixa de marcação (checkbox) com rótulo
% --------------------------------------------------------------
\caixa{exemplo_opcao1}{Opção 1}
\hspace{1cm}
\caixa{exemplo_opcao2}{Opção 2}
\hspace{1cm}
\caixa{exemplo_opcao3}{Opção 3}

\vspace{0.5cm}

% --------------------------------------------------------------
% 4. e 5. Campo de texto e checkbox dentro de tabela
% --------------------------------------------------------------
% Use coluna "m{...}" (em vez de "p{...}") para os campos ficarem
% centralizados também na altura da célula.
\begin{center}
\begin{tabular}{|>{\centering\arraybackslash}m{4cm}|>{\centering\arraybackslash}m{3cm}|>{\centering\arraybackslash}m{2cm}|}
    \hline
    \textbf{Item} & \textbf{Valor} & \textbf{OK?} \\
    \hline
    Item 1 & \campotabela{exemplo_tab_item1}{1cm} & \checktabela{exemplo_tab_check1} \\
    \hline
    Item 2 & \campotabela{exemplo_tab_item2}{1cm} & \checktabela{exemplo_tab_check2} \\
    \hline
\end{tabular}
\end{center}

\end{Form} dentro de main.tex (onde hoje
% está \input{secoes/teste.tex}).
%
% As macros usadas abaixo (\campotexto, \campolongo, \caixa,
% \campotabela, \checktabela) são definidas em
% estrut/campos-formulario.tex — não precisa mexer lá, só usá-las.
%
% REGRA OBRIGATÓRIA: todo campo só funciona dentro do ambiente Form,
% como abaixo. NOMES (1º argumento de cada macro) DEVEM SER ÚNICOS
% no documento inteiro.
% ====================================================================

\begin{Form}

\section{Exemplo de Formulário}

% --------------------------------------------------------------
% 1. Campo de texto simples (uma linha)
% --------------------------------------------------------------
Nome do responsável: \campotexto{exemplo_nome}{6cm}

\vspace{0.5cm}

% --------------------------------------------------------------
% 2. Campo de texto longo (múltiplas linhas)
% --------------------------------------------------------------
Observações gerais:

\campolongo{exemplo_observacoes}{3cm}

\vspace{0.5cm}

% --------------------------------------------------------------
% 3. Caixa de marcação (checkbox) com rótulo
% --------------------------------------------------------------
\caixa{exemplo_opcao1}{Opção 1}
\hspace{1cm}
\caixa{exemplo_opcao2}{Opção 2}
\hspace{1cm}
\caixa{exemplo_opcao3}{Opção 3}

\vspace{0.5cm}

% --------------------------------------------------------------
% 4. e 5. Campo de texto e checkbox dentro de tabela
% --------------------------------------------------------------
% Use coluna "m{...}" (em vez de "p{...}") para os campos ficarem
% centralizados também na altura da célula.
\begin{center}
\begin{tabular}{|>{\centering\arraybackslash}m{4cm}|>{\centering\arraybackslash}m{3cm}|>{\centering\arraybackslash}m{2cm}|}
    \hline
    \textbf{Item} & \textbf{Valor} & \textbf{OK?} \\
    \hline
    Item 1 & \campotabela{exemplo_tab_item1}{1cm} & \checktabela{exemplo_tab_check1} \\
    \hline
    Item 2 & \campotabela{exemplo_tab_item2}{1cm} & \checktabela{exemplo_tab_check2} \\
    \hline
\end{tabular}
\end{center}

\end{Form} dentro de main.tex (onde hoje
% está \input{secoes/teste.tex}).
%
% As macros usadas abaixo (\campotexto, \campolongo, \caixa,
% \campotabela, \checktabela) são definidas em
% estrut/campos-formulario.tex — não precisa mexer lá, só usá-las.
%
% REGRA OBRIGATÓRIA: todo campo só funciona dentro do ambiente Form,
% como abaixo. NOMES (1º argumento de cada macro) DEVEM SER ÚNICOS
% no documento inteiro.
% ====================================================================

\begin{Form}

\section{Exemplo de Formulário}

% --------------------------------------------------------------
% 1. Campo de texto simples (uma linha)
% --------------------------------------------------------------
Nome do responsável: \campotexto{exemplo_nome}{6cm}

\vspace{0.5cm}

% --------------------------------------------------------------
% 2. Campo de texto longo (múltiplas linhas)
% --------------------------------------------------------------
Observações gerais:

\campolongo{exemplo_observacoes}{3cm}

\vspace{0.5cm}

% --------------------------------------------------------------
% 3. Caixa de marcação (checkbox) com rótulo
% --------------------------------------------------------------
\caixa{exemplo_opcao1}{Opção 1}
\hspace{1cm}
\caixa{exemplo_opcao2}{Opção 2}
\hspace{1cm}
\caixa{exemplo_opcao3}{Opção 3}

\vspace{0.5cm}

% --------------------------------------------------------------
% 4. e 5. Campo de texto e checkbox dentro de tabela
% --------------------------------------------------------------
% Use coluna "m{...}" (em vez de "p{...}") para os campos ficarem
% centralizados também na altura da célula.
\begin{center}
\begin{tabular}{|>{\centering\arraybackslash}m{4cm}|>{\centering\arraybackslash}m{3cm}|>{\centering\arraybackslash}m{2cm}|}
    \hline
    \textbf{Item} & \textbf{Valor} & \textbf{OK?} \\
    \hline
    Item 1 & \campotabela{exemplo_tab_item1}{1cm} & \checktabela{exemplo_tab_check1} \\
    \hline
    Item 2 & \campotabela{exemplo_tab_item2}{1cm} & \checktabela{exemplo_tab_check2} \\
    \hline
\end{tabular}
\end{center}

\end{Form} dentro de main.tex (onde hoje
% está \input{secoes/teste.tex}).
%
% As macros usadas abaixo (\campotexto, \campolongo, \caixa,
% \campotabela, \checktabela) são definidas em
% estrut/campos-formulario.tex — não precisa mexer lá, só usá-las.
%
% REGRA OBRIGATÓRIA: todo campo só funciona dentro do ambiente Form,
% como abaixo. NOMES (1º argumento de cada macro) DEVEM SER ÚNICOS
% no documento inteiro.
% ====================================================================

\begin{Form}

\section{Exemplo de Formulário}

% --------------------------------------------------------------
% 1. Campo de texto simples (uma linha)
% --------------------------------------------------------------
Nome do responsável: \campotexto{exemplo_nome}{6cm}

\vspace{0.5cm}

% --------------------------------------------------------------
% 2. Campo de texto longo (múltiplas linhas)
% --------------------------------------------------------------
Observações gerais:

\campolongo{exemplo_observacoes}{3cm}

\vspace{0.5cm}

% --------------------------------------------------------------
% 3. Caixa de marcação (checkbox) com rótulo
% --------------------------------------------------------------
\caixa{exemplo_opcao1}{Opção 1}
\hspace{1cm}
\caixa{exemplo_opcao2}{Opção 2}
\hspace{1cm}
\caixa{exemplo_opcao3}{Opção 3}

\vspace{0.5cm}

% --------------------------------------------------------------
% 4. e 5. Campo de texto e checkbox dentro de tabela
% --------------------------------------------------------------
% Use coluna "m{...}" (em vez de "p{...}") para os campos ficarem
% centralizados também na altura da célula.
\begin{center}
\begin{tabular}{|>{\centering\arraybackslash}m{4cm}|>{\centering\arraybackslash}m{3cm}|>{\centering\arraybackslash}m{2cm}|}
    \hline
    \textbf{Item} & \textbf{Valor} & \textbf{OK?} \\
    \hline
    Item 1 & \campotabela{exemplo_tab_item1}{1cm} & \checktabela{exemplo_tab_check1} \\
    \hline
    Item 2 & \campotabela{exemplo_tab_item2}{1cm} & \checktabela{exemplo_tab_check2} \\
    \hline
\end{tabular}
\end{center}

\end{Form}